\section{Quadratic curvature actions and coupled fields}
\label{sec:quadratic-matter}

The scalar-free estimate and the traceless-Ricci window have direct
consequences for curvature actions and Euclidean coupled equations.  We
give the normalizations and algebra explicitly.  These consequences do
not require the existence of a minimizing metric or of a solution to a
coupled field equation.

Let $M$ be a closed connected oriented smooth four-manifold.  Write
\[
 V=\Vol_g(M),\quad S=\int_Ms^2\,d\mu_g,\quad
 P=\int_M|\Ric^\circ|^2\,d\mu_g,\quad
 R=\int_M|\Ric|^2\,d\mu_g=P+\frac S4,
 \quad W_\pm=\int_M|W^\pm|^2\,d\mu_g.
\]
Ricci curvature has the tensor norm, two-forms have the exterior-algebra
norm, and Weyl curvature has the Hilbert--Schmidt norm as an operator on
two-forms.  Set $T_\pm=2\chi(M)\pm3\tau(M)$.

For the estimates involving topology, assume that $A>0$ is fixed and
\begin{equation}\label{eq:qm-supply}
 T_+=A(1-h),\qquad 0<h\le\frac13.
\end{equation}
For every metric under consideration, assume there is a Spin$^c$
structure whose unperturbed Seiberg--Witten equations are solvable for
every metric in its conformal class, and whose real first Chern class
$a$ satisfies
\[
 p=(a_g^+)^2\ge A,\qquad a^2\ge T_+.
\]
The selected structure may depend on $g$, but is fixed throughout its
conformal class.  These are the actual input hypotheses of
Theorems~\ref{thm:lc-scalar-free} and~\ref{thm:linear-window}; a
formal characteristic class with the same square is not sufficient.
Define
\[
 D_2=\frac32S+18W_+-72\pi^2p,\qquad
 \Sigma=S-32\pi^2p.
\]
Both quantities are nonnegative.  Theorem~\ref{thm:lc-scalar-free}
supplies a universal constant $C>0$ such that
\begin{equation}\label{eq:qm-slack}
 \Sigma\le C(P+D_2).
\end{equation}

\subsection{A quadratic-action gap and an absolute Ricci budget}

\begin{proposition}[Positive-coupling curvature actions]
\label{prop:qm-action}
Let $u,v>0$ be fixed and put
\[
 \mathscr I_{u,v}(g)=uP+v\left(\frac32S+18W_+\right),
 \qquad K_C=\frac92+\frac98C.
\]
Under the stated supply hypotheses,
\begin{equation}\label{eq:qm-action-gap}
 \mathscr I_{u,v}(g)-72\pi^2vA
 \ge\frac{36\pi^2\min\{u,v\}}{K_C}\,Ah.
\end{equation}
In terms of the full Weyl operator, the same action is
\begin{equation}\label{eq:qm-action-signature}
 \mathscr I_{u,v}(g)
 =uP+\frac32vS+9v\int_M|W|^2\,d\mu_g+108\pi^2v\tau(M).
\end{equation}
The proof of \eqref{eq:qm-action-gap} only needs $h>0$, not the upper
bound in \eqref{eq:qm-supply}.
\end{proposition}

\begin{proof}
The chiral Chern--Weil identity is
\[
 4\pi^2T_+=\frac S{24}+2W_+-\frac P2.
\]
Substituting it into the definitions gives two exact identities:
\begin{equation}\label{eq:qm-exact-identities}
 D_2=\frac92R-36\pi^2(2p-T_+),\qquad
 36\pi^2(p-T_+)=\frac92P+\frac98\Sigma-D_2.
\end{equation}
Since $p-T_+\ge Ah$, \eqref{eq:qm-slack} and $D_2\ge0$ imply
\[
 36\pi^2Ah
 \le\frac92P+\frac98\Sigma-D_2
 \le K_C(P+D_2).
\]
On the other hand,
\[
 \mathscr I_{u,v}-72\pi^2vA
 =uP+vD_2+72\pi^2v(p-A)
 \ge\min\{u,v\}(P+D_2).
\]
This proves the gap.  Finally the signature identity
$W_+-W_-=12\pi^2\tau(M)$ gives
$W_+=\frac12(W_++W_-)+6\pi^2\tau(M)$, proving
\eqref{eq:qm-action-signature}.
\end{proof}

The subtracted budget in \eqref{eq:qm-action-gap} uses the fixed number
$A$, not the metric-dependent projection $p$.  No assertion about
sharpness or attainment of this action bound is implicit in the estimate.

For the following consequences set
\begin{equation}\label{eq:qm-mu}
 U=\frac{3\min\{T_+,T_-\}}A,\qquad L_h=1+(1+c_0)h,
 \qquad
 \mu=\begin{cases}
 \displaystyle\frac{(L_h-U)_+}{L_h+U/18},&U>0,\\[6pt]
 1,&U\le0,
 \end{cases}
\end{equation}
where $c_0>0$ is the universal constant in
Theorem~\ref{thm:linear-window}.  Thus $0\le\mu\le1$ and that theorem
gives $\J(g)\ge\mu$ for every metric with the stated supply.  Put
\begin{equation}\label{eq:qm-P-bound}
 \mathscr P(A,h,U)=2\pi^2A(1+h)\mu.
\end{equation}
The first identity in \eqref{eq:qm-exact-identities} and $D_2\ge0$ give
\[
 R\ge8\pi^2(2p-T_+)\ge8\pi^2A(1+h).
\]
Since $P=R\J/4$, we have proved the absolute bound
\begin{equation}\label{eq:qm-absolute-P}
 P\ge\mathscr P(A,h,U).
\end{equation}
This bound is positive whenever $U<L_h$.  On the balanced families of
Theorem~\ref{thm:vg-balanced-family},
$U_j=1+h_j+O(h_j^2)$ and $h_j\to0$.  Formula~\eqref{eq:qm-mu}
then gives $\mu_j/h_j\to18c_0/19$.  Consequently there are constants
$c_{\mathrm J},C_0>0$, independent of $j$, such that, for all sufficiently
large $j$,
\begin{equation}\label{eq:qm-balanced-P}
 \J(g)\ge c_{\mathrm J}h_j,\qquad
 P(g)\ge C_0A_jh_j\quad\hbox{for every smooth }g,
 \qquad C_0=2\pi^2c_{\mathrm J}.
\end{equation}

\subsection{The mixed chiral curvature block}

\begin{proposition}\label{prop:qm-chiral-block}
Relative to $\Lambda^2=\Lambda^+\oplus\Lambda^-$, write the Riemann
curvature operator as
\[
 \mathcal R_g=
 \begin{pmatrix}
 W^++\frac{s}{12}\mathrm{Id}&B\\
 B^*&W^-+\frac{s}{12}\mathrm{Id}
 \end{pmatrix}.
\]
Then
\begin{equation}\label{eq:qm-block-norm}
 |B|_{\mathrm{HS}}^2=\frac14|\Ric^\circ|^2,
 \qquad 4\int_M|B|_{\mathrm{HS}}^2\,d\mu_g=P.
\end{equation}
In particular the integral on the left of the second identity is at
least $\mathscr P(A,h,U)$ under the supply hypotheses.
\end{proposition}

\begin{proof}
At a point diagonalize $\Ric^\circ$ with eigenvalues $q_1,\ldots,q_4$,
where $\sum_iq_i=0$.  Use the orthonormal chiral bases
$(e^{1j}\pm *e^{1j})/\sqrt2$, $j=2,3,4$.  The traceless-Ricci part
of the curvature operator has mixed block with diagonal entries,
up to signs,
\[
 \frac{q_1+q_2}{2},\qquad
 \frac{q_1+q_3}{2},\qquad
 \frac{q_1+q_4}{2}.
\]
Their squared sum is
\[
 \frac14\left(3q_1^2+2q_1\sum_{i=2}^4q_i+
                          \sum_{i=2}^4q_i^2\right)
 =\frac14\sum_{i=1}^4q_i^2.
\]
This proves the pointwise identity; integration and
\eqref{eq:qm-absolute-P} prove the rest.
\end{proof}

Thus the bound measures a compulsory failure of the Levi-Civita curvature
operator to preserve the two chiral summands.  It concerns the curvature
of one common metric, not arbitrary independently chosen gauge
connections.  The operator norm in \eqref{eq:qm-block-norm} fixes the
representation constant and should not be replaced by an unspecified
Lie-algebra trace norm.

\subsection{Trace-free stress and Maxwell or Yang--Mills overlap}

\begin{proposition}[Conditional estimates for coupled solutions]
\label{prop:qm-stress}
Suppose a Riemannian metric with the stated supply is part of a smooth
solution of
\begin{equation}\label{eq:qm-einstein-matter}
 \Ric-\frac12s\,g+\Lambda g=\kappa T,
 \qquad \Lambda\in\mathbb R,\quad \kappa\ne0,
\end{equation}
where $T$ is the stress tensor in the specified normalization.  Then
\begin{equation}\label{eq:qm-stress-bound}
 \int_M|T^\circ|^2\,d\mu_g
 =\frac P{\kappa^2}\ge\frac{\mathscr P(A,h,U)}{\kappa^2}.
\end{equation}
If $T$ is trace-free, necessarily $s=4\Lambda$ and $\Lambda<0$, and
\begin{equation}\label{eq:qm-stress-dimensionless}
 \frac{P}{\Lambda^2V}\ge\frac{4\mu}{4-\mu}.
\end{equation}
On the balanced families this also gives, for all sufficiently large $j$,
$\int|T^\circ|^2\ge C_0A_jh_j/\kappa^2$ and
$P/(\Lambda^2V)\ge4c_{\mathrm J}h_j/(4-c_{\mathrm J}h_j)$.
\end{proposition}

\begin{proof}
The trace-free part of \eqref{eq:qm-einstein-matter} is
$\Ric^\circ=\kappa T^\circ$, which proves
\eqref{eq:qm-stress-bound}.  If $\tr_gT=0$, the trace equation is
$-s+4\Lambda=0$.  The actual scalar inequality
$\int s_-^2\,d\mu_g\ge32\pi^2p>0$ rules out nonnegative constant
scalar curvature, so $\Lambda<0$.  Now $R=P+4\Lambda^2V$ and
\[
 \frac{4P}{P+4\Lambda^2V}=\J(g)\ge\mu.
\]
Rearranging proves \eqref{eq:qm-stress-dimensionless}, since
$\mu\le1<4$.  Use \eqref{eq:qm-balanced-P} for the final assertions.
\end{proof}

For a real Maxwell two-form use the normalization
\begin{equation}\label{eq:qm-maxwell-stress}
 T_{ij}=F_{ik}F_j{}^k-\frac14g_{ij}F_{k\ell}F^{k\ell},
 \qquad |F|^2=\frac12F_{ij}F^{ij}.
\end{equation}
An Einstein--Maxwell solution also satisfies $dF=d{*F}=0$.
For Yang--Mills curvature, choose a positive invariant inner product on
the gauge Lie algebra, sum \eqref{eq:qm-maxwell-stress} over an
orthonormal Lie-algebra basis, and impose the corresponding Yang--Mills
equation.  The following estimate needs only the resulting stress tensor
and \eqref{eq:qm-einstein-matter}.

\begin{corollary}[Chiral overlap]\label{cor:qm-overlap}
Every such Einstein--Maxwell or Einstein--Yang--Mills solution with the
stated supply satisfies
\begin{equation}\label{eq:qm-overlap}
 \int_M|F^+|^2|F^-|^2\,d\mu_g
 \ge\frac{\mathscr P(A,h,U)}{4\kappa^2}.
\end{equation}
For Maxwell fields the exact identity behind this estimate is
\begin{equation}\label{eq:qm-maxwell-identity}
 |T|^2=4|F^+|^2|F^-|^2.
\end{equation}
For Lie-algebra-valued curvature one has instead
$|T|\le2|F^+||F^-|$.
\end{corollary}

\begin{proof}
At a point choose an oriented orthonormal frame with
$F=a\,e^{12}+b\,e^{34}$.  Then
\[
 |F^+|^2=\frac{(a+b)^2}{2},\qquad
 |F^-|^2=\frac{(a-b)^2}{2}.
\]
The diagonal entries of $T$ are two copies of $(a^2-b^2)/2$ and two
copies of its negative.  Thus
$|T|^2=(a^2-b^2)^2=4|F^+|^2|F^-|^2$.
For a Lie-algebra-valued form write $F=\sum_\alpha F^\alpha t_\alpha$
in an orthonormal basis.  By the real-form identity, the triangle
inequality, and Cauchy--Schwarz,
\[
 |T|\le\sum_\alpha|T(F^\alpha)|
 =2\sum_\alpha|(F^\alpha)^+||(F^\alpha)^-|
 \le2|F^+||F^-|.
\]
In both cases the stress tensor is trace-free.  Integrating the squared
inequality and applying Proposition~\ref{prop:qm-stress} proves
\eqref{eq:qm-overlap}.
\end{proof}

When $\mu>0$, no solution in this class has a purely self-dual or purely
anti-self-dual real field.  The overlap estimate does not separately
bound the two $L^2$ energies without additional information: one factor
can grow or concentrate while the other decreases.  The assertions are
conditional on an actual coupled solution, and neither construct such
a solution nor determine the complete set of admissible fluxes.

\subsection{The leading Ricci residual modulo constant rescaling}

\begin{proposition}\label{prop:qm-ricci-residual}
For any constant $\gamma\in\mathbb R$, define
$\mathcal B_\gamma(g)=\gamma\Ric_g$.  If
$I=\int_Ms\,d\mu_g$ and $L_s=S-I^2/V$, then
\begin{equation}\label{eq:qm-residual}
 \inf_{\lambda\in\mathbb R}
   \int_M|\mathcal B_\gamma(g)-\lambda g|^2\,d\mu_g
 =\gamma^2\left(P+\frac{L_s}{4}\right)
 =\frac{\gamma^2R}{4}\E(g)
 \ge\gamma^2\mathscr P(A,h,U).
\end{equation}
The infimum is attained at $\lambda=\gamma I/(4V)$.
\end{proposition}

\begin{proof}
The trace-free and pure-trace summands are orthogonal, so
\[
 \int_M|\gamma\Ric-\lambda g|^2\,d\mu_g
 =\gamma^2P+4\int_M\left(\frac{\gamma s}{4}-\lambda\right)^2\,d\mu_g.
\]
The last term is minimized by its mean and has minimum
$\gamma^2L_s/4$.  The identity
$\E=(4P+L_s)/R$ and \eqref{eq:qm-absolute-P} give the remaining claims.
\end{proof}

In a model normalized so that its leading metric beta tensor is
$\beta^{(1)}=\alpha'\Ric$, this is the residual obtained by removing
only the constant scaling direction, with $\gamma=\alpha'$.
It is not a bound after minimizing over Lie derivatives
$\mathcal L_Xg$, additional background fields, or higher-order terms.
All estimates in this section use a positive-definite metric and the
real splitting $*^2=1$ on two-forms; no Lorentzian continuation is
asserted.
